% =====================================================================
%  Anatomía de un LLM — Fase I — Módulo 4
%  Cálculo para el aprendizaje automático
%  Compilar: latexmk -r ../latexmkrc M04-calculo.tex
% =====================================================================
\documentclass[11pt,a4paper,twoside,openany]{book}
\usepackage{../estilo-apuntes}

\newcommand{\ApuntesCurso}{Anatomía de un LLM}
\newcommand{\ApuntesModulo}{Módulo 4}

\hypersetup{pdftitle={Módulo 4. Cálculo para el aprendizaje automático (resumen)},pdfauthor={Anatomía de un LLM}}

% Notación propia del módulo
\newcommand{\diag}{\operatorname{diag}}
\newcommand{\tr}{\operatorname{tr}}
\newcommand{\Var}{\operatorname{Var}}
\newcommand{\dom}{\operatorname{dom}}
\newcommand{\T}{^{\top}}
\newcommand{\uno}{\mathbf{1}}

\makeatletter
\newlabel{ej:parciales}{{1.2}{1}{}{ejemplo.2}{}}
\newlabel{ej:parciales@cref}{{[ejemplo][2][]1.2}{[1][1][]1}}
\newlabel{ej:direccional}{{1.10}{2}{}{ejemplo.10}{}}
\newlabel{ej:direccional@cref}{{[ejemplo][10][]1.10}{[1][2][]2}}
\newlabel{fig:nivel}{{1.1}{4}{Curvas de nivel $x^{2}+3y^{2}=c$ para $c=1$; $2{,}5$; $4$. En $(1,1)$ el gradiente (flecha, a escala reducida) es perpendicular a la tangente a la curva $f=4$ (trazo discontinuo), de dirección $(3,-1)\T $}{figure.1.1}{}}
\newlabel{fig:nivel@cref}{{[figure][1][1]1.1}{[1][3][]4}}
\newlabel{ej:dg}{{1.16}{4}{}{ejemplo.16}{}}
\newlabel{ej:dg@cref}{{[ejemplo][16][]1.16}{[1][4][]4}}
\newlabel{eje:1-vf}{{1.1}{5}{}{ejercicio.1.1}{}}
\newlabel{eje:1-vf@cref}{{[ejercicio][1][1]1.1}{[1][4][]5}}
\newlabel{eje:1-gradiente}{{1.2}{5}{}{ejercicio.1.2}{}}
\newlabel{eje:1-gradiente@cref}{{[ejercicio][2][1]1.2}{[1][5][]5}}
\newlabel{eje:1-d2l}{{1.3}{5}{}{ejercicio.1.3}{}}
\newlabel{eje:1-d2l@cref}{{[ejercicio][3][1]1.3}{[1][5][]5}}
\newlabel{eje:1-norma}{{1.4}{5}{}{ejercicio.1.4}{}}
\newlabel{eje:1-norma@cref}{{[ejercicio][4][1]1.4}{[1][5][]5}}
\newlabel{eje:1-mit}{{1.5}{5}{}{ejercicio.1.5}{}}
\newlabel{eje:1-mit@cref}{{[ejercicio][5][1]1.5}{[1][5][]5}}
\newlabel{eje:1-minimo}{{1.6}{5}{}{ejercicio.1.6}{}}
\newlabel{eje:1-minimo@cref}{{[ejercicio][6][1]1.6}{[1][5][]5}}
\newlabel{ej:jac-lineal}{{2.2}{6}{}{ejemplo.2}{}}
\newlabel{ej:jac-lineal@cref}{{[ejemplo][2][]2.2}{[1][6][]6}}
\newlabel{ej:jac}{{2.3}{6}{}{ejemplo.3}{}}
\newlabel{ej:jac@cref}{{[ejemplo][3][]2.3}{[1][6][]6}}
\newlabel{ej:cadena}{{2.6}{7}{}{ejemplo.6}{}}
\newlabel{ej:cadena@cref}{{[ejemplo][6][]2.6}{[1][7][]7}}
\newlabel{fig:grafo}{{2.1}{8}{Grafo computacional del \cref {ej:cadena}. Dentro de cada nodo figura el valor calculado en el paso hacia delante; junto a él, en gris, el adjunto calculado en el paso hacia atrás con la \cref {prop:retro}. La entrada $x$ recibe contribuciones por dos caminos, que se suman}{figure.2.1}{}}
\newlabel{fig:grafo@cref}{{[figure][1][2]2.1}{[1][7][]8}}
\newlabel{eje:2-formas}{{2.1}{9}{}{ejercicio.2.1}{}}
\newlabel{eje:2-formas@cref}{{[ejercicio][1][2]2.1}{[1][9][]9}}
\newlabel{eje:2-producto}{{2.2}{9}{}{ejercicio.2.2}{}}
\newlabel{eje:2-producto@cref}{{[ejercicio][2][2]2.2}{[1][9][]9}}
\newlabel{eje:2-jacobiano}{{2.3}{9}{}{ejercicio.2.3}{}}
\newlabel{eje:2-jacobiano@cref}{{[ejercicio][3][2]2.3}{[1][9][]9}}
\newlabel{eje:2-d2l}{{2.4}{9}{}{ejercicio.2.4}{}}
\newlabel{eje:2-d2l@cref}{{[ejercicio][4][2]2.4}{[1][9][]9}}
\newlabel{eje:2-coste}{{2.5}{9}{}{ejercicio.2.5}{}}
\newlabel{eje:2-coste@cref}{{[ejercicio][5][2]2.5}{[1][9][]9}}
\newlabel{eje:2-escalar}{{2.6}{9}{}{ejercicio.2.6}{}}
\newlabel{eje:2-escalar@cref}{{[ejercicio][6][2]2.6}{[1][9][]9}}
\newlabel{ej:capa}{{3.4}{11}{}{ejemplo.4}{}}
\newlabel{ej:capa@cref}{{[ejemplo][4][]3.4}{[1][11][]11}}
\newlabel{ej:p-y}{{3.9}{12}{}{ejemplo.9}{}}
\newlabel{ej:p-y@cref}{{[ejemplo][9][]3.9}{[1][12][]12}}
\newlabel{eje:3-embeddings}{{3.1}{13}{}{ejercicio.3.1}{}}
\newlabel{eje:3-embeddings@cref}{{[ejercicio][1][3]3.1}{[1][13][]13}}
\newlabel{eje:3-signos}{{3.2}{13}{}{ejercicio.3.2}{}}
\newlabel{eje:3-signos@cref}{{[ejercicio][2][3]3.2}{[1][13][]13}}
\newlabel{eje:3-capa}{{3.3}{13}{}{ejercicio.3.3}{}}
\newlabel{eje:3-capa@cref}{{[ejercicio][3][3]3.3}{[1][13][]13}}
\newlabel{eje:3-softmax}{{3.4}{13}{}{ejercicio.3.4}{}}
\newlabel{eje:3-softmax@cref}{{[ejercicio][4][3]3.4}{[1][13][]13}}
\newlabel{eje:3-nucleo}{{3.5}{13}{}{ejercicio.3.5}{}}
\newlabel{eje:3-nucleo@cref}{{[ejercicio][5][3]3.5}{[1][13][]13}}
\newlabel{eje:3-word2vec}{{3.6}{13}{}{ejercicio.3.6}{}}
\newlabel{eje:3-word2vec@cref}{{[ejercicio][6][3]3.6}{[1][13][]13}}
\newlabel{ej:hessiano}{{4.3}{14}{}{ejemplo.3}{}}
\newlabel{ej:hessiano@cref}{{[ejemplo][3][]4.3}{[1][14][]14}}
\newlabel{ej:taylor}{{4.6}{15}{}{ejemplo.6}{}}
\newlabel{ej:taylor@cref}{{[ejemplo][6][]4.6}{[1][15][]15}}
\newlabel{ej:dg-cuad}{{4.13}{16}{}{ejemplo.13}{}}
\newlabel{ej:dg-cuad@cref}{{[ejemplo][13][]4.13}{[1][16][]16}}
\newlabel{eje:4-vf}{{4.1}{17}{}{ejercicio.4.1}{}}
\newlabel{eje:4-vf@cref}{{[ejercicio][1][4]4.1}{[1][17][]17}}
\newlabel{eje:4-taylor}{{4.2}{17}{}{ejercicio.4.2}{}}
\newlabel{eje:4-taylor@cref}{{[ejercicio][2][4]4.2}{[1][17][]17}}
\newlabel{eje:4-mit}{{4.3}{18}{}{ejercicio.4.3}{}}
\newlabel{eje:4-mit@cref}{{[ejercicio][3][4]4.3}{[1][18][]18}}
\newlabel{eje:4-tasa}{{4.4}{18}{}{ejercicio.4.4}{}}
\newlabel{eje:4-tasa@cref}{{[ejercicio][4][4]4.4}{[1][18][]18}}
\newlabel{eje:4-necesaria}{{4.5}{18}{}{ejercicio.4.5}{}}
\newlabel{eje:4-necesaria@cref}{{[ejercicio][5][4]4.5}{[1][18][]18}}
\newlabel{fig:convexa}{{5.1}{19}{Una función convexa queda por debajo de sus cuerdas (\cref {def:convexa}) y por encima de sus tangentes (\cref {teo:primer-orden})}{figure.5.1}{}}
\newlabel{fig:convexa@cref}{{[figure][1][5]5.1}{[1][19][]19}}
\newlabel{ej:cuad-convexa}{{5.6}{20}{}{ejemplo.6}{}}
\newlabel{ej:cuad-convexa@cref}{{[ejemplo][6][]5.6}{[1][20][]20}}
\newlabel{ej:red-lineal}{{5.13}{22}{Una red lineal de dos capas}{ejemplo.13}{}}
\newlabel{ej:red-lineal@cref}{{[ejemplo][13][]5.13}{[1][22][]22}}
\newlabel{fig:hiperbola}{{5.2}{23}{Mínimos globales de $L(w_1,w_2)=(1-w_1w_2)^{2}$ (\cref {ej:red-lineal}). El segmento que une dos mínimos sale del conjunto de mínimos: la pérdida no es convexa}{figure.5.2}{}}
\newlabel{fig:hiperbola@cref}{{[figure][2][5]5.2}{[1][22][]23}}
\newlabel{eje:5-vf}{{5.1}{23}{}{ejercicio.5.1}{}}
\newlabel{eje:5-vf@cref}{{[ejercicio][1][5]5.1}{[1][23][]23}}
\newlabel{eje:5-hessianos}{{5.2}{23}{}{ejercicio.5.2}{}}
\newlabel{eje:5-hessianos@cref}{{[ejercicio][2][5]5.2}{[1][23][]23}}
\newlabel{eje:5-boyd}{{5.3}{23}{}{ejercicio.5.3}{}}
\newlabel{eje:5-boyd@cref}{{[ejercicio][3][5]5.3}{[1][23][]23}}
\newlabel{eje:5-red}{{5.4}{23}{}{ejercicio.5.4}{}}
\newlabel{eje:5-red@cref}{{[ejercicio][4][5]5.4}{[1][23][]23}}
\newlabel{eje:5-softmax-W}{{5.5}{23}{}{ejercicio.5.5}{}}
\newlabel{eje:5-softmax-W@cref}{{[ejercicio][5][5]5.5}{[1][23][]23}}
\newlabel{eje:5-unico}{{5.6}{23}{}{ejercicio.5.6}{}}
\newlabel{eje:5-unico@cref}{{[ejercicio][6][5]5.6}{[1][23][]23}}
\newlabel{ej:exp}{{6.3}{25}{}{ejemplo.3}{}}
\newlabel{ej:exp@cref}{{[ejemplo][3][]6.3}{[1][25][]25}}
\newlabel{fig:error-h}{{6.1}{27}{Error de la diferencia centrada para $f(x)=e^{x}$ en $x=0$, calculada en doble precisión (puntos), y cota $E(h)$ de la \cref {prop:h-optimo} con $\varepsilon =1{,}1\cdot 10^{-16}$ y $\abs {f}=M=1$ (línea). A la derecha domina el truncamiento (pendiente~2); a la izquierda, el redondeo (pendiente~$-1$). El error medido mínimo, $1{,}2\cdot 10^{-11}$, se obtiene con $h=10^{-5}$}{figure.6.1}{}}
\newlabel{fig:error-h@cref}{{[figure][1][6]6.1}{[1][26][]27}}
\newlabel{eje:6-porque}{{6.1}{28}{}{ejercicio.6.1}{}}
\newlabel{eje:6-porque@cref}{{[ejercicio][1][6]6.1}{[1][28][]28}}
\newlabel{eje:6-cubo}{{6.2}{28}{}{ejercicio.6.2}{}}
\newlabel{eje:6-cubo@cref}{{[ejercicio][2][6]6.2}{[1][28][]28}}
\newlabel{eje:6-direccional}{{6.3}{28}{}{ejercicio.6.3}{}}
\newlabel{eje:6-direccional@cref}{{[ejercicio][3][6]6.3}{[1][28][]28}}
\newlabel{eje:6-cs231n}{{6.4}{28}{}{ejercicio.6.4}{}}
\newlabel{eje:6-cs231n@cref}{{[ejercicio][4][6]6.4}{[1][28][]28}}
\newlabel{eje:6-fermi}{{6.5}{28}{}{ejercicio.6.5}{}}
\newlabel{eje:6-fermi@cref}{{[ejercicio][5][6]6.5}{[1][28][]28}}
\newlabel{eje:6-hoptimo}{{6.6}{28}{}{ejercicio.6.6}{}}
\newlabel{eje:6-hoptimo@cref}{{[ejercicio][6][6]6.6}{[1][28][]28}}
\newlabel{ap:soluciones}{{A}{29}{Soluciones de los ejercicios}{appendix.A}{}}
\newlabel{ap:soluciones@cref}{{[appendix][1][2147483647]A}{[1][29][]29}}
\newlabel{sol:eje:1-vf}{{A.1}{29}{Tema 1}{section*.3}{}}
\newlabel{sol:eje:1-vf@cref}{{[subappendix][1][2147483647,1]A.1}{[1][29][]29}}
\newlabel{sol:eje:1-gradiente}{{A.1}{29}{Tema 1}{section*.4}{}}
\newlabel{sol:eje:1-gradiente@cref}{{[subappendix][1][2147483647,1]A.1}{[1][29][]29}}
\newlabel{sol:eje:1-d2l}{{A.1}{30}{Tema 1}{section*.5}{}}
\newlabel{sol:eje:1-d2l@cref}{{[subappendix][1][2147483647,1]A.1}{[1][29][]30}}
\newlabel{sol:eje:1-norma}{{A.1}{30}{Tema 1}{section*.6}{}}
\newlabel{sol:eje:1-norma@cref}{{[subappendix][1][2147483647,1]A.1}{[1][30][]30}}
\newlabel{sol:eje:1-mit}{{A.1}{30}{Tema 1}{section*.7}{}}
\newlabel{sol:eje:1-mit@cref}{{[subappendix][1][2147483647,1]A.1}{[1][30][]30}}
\newlabel{sol:eje:1-minimo}{{A.1}{31}{Tema 1}{section*.8}{}}
\newlabel{sol:eje:1-minimo@cref}{{[subappendix][1][2147483647,1]A.1}{[1][31][]31}}
\newlabel{sol:eje:2-formas}{{A.2}{31}{Tema 2}{section*.9}{}}
\newlabel{sol:eje:2-formas@cref}{{[subappendix][2][2147483647,1]A.2}{[1][31][]31}}
\newlabel{sol:eje:2-producto}{{A.2}{32}{Tema 2}{section*.10}{}}
\newlabel{sol:eje:2-producto@cref}{{[subappendix][2][2147483647,1]A.2}{[1][32][]32}}
\newlabel{sol:eje:2-jacobiano}{{A.2}{32}{Tema 2}{section*.11}{}}
\newlabel{sol:eje:2-jacobiano@cref}{{[subappendix][2][2147483647,1]A.2}{[1][32][]32}}
\newlabel{sol:eje:2-d2l}{{A.2}{32}{Tema 2}{section*.12}{}}
\newlabel{sol:eje:2-d2l@cref}{{[subappendix][2][2147483647,1]A.2}{[1][32][]32}}
\newlabel{sol:eje:2-coste}{{A.2}{33}{Tema 2}{section*.13}{}}
\newlabel{sol:eje:2-coste@cref}{{[subappendix][2][2147483647,1]A.2}{[1][33][]33}}
\newlabel{sol:eje:2-escalar}{{A.2}{33}{Tema 2}{section*.14}{}}
\newlabel{sol:eje:2-escalar@cref}{{[subappendix][2][2147483647,1]A.2}{[1][33][]33}}
\newlabel{sol:eje:3-embeddings}{{A.3}{34}{Tema 3}{section*.15}{}}
\newlabel{sol:eje:3-embeddings@cref}{{[subappendix][3][2147483647,1]A.3}{[1][33][]34}}
\newlabel{sol:eje:3-signos}{{A.3}{34}{Tema 3}{section*.16}{}}
\newlabel{sol:eje:3-signos@cref}{{[subappendix][3][2147483647,1]A.3}{[1][34][]34}}
\newlabel{sol:eje:3-capa}{{A.3}{34}{Tema 3}{section*.17}{}}
\newlabel{sol:eje:3-capa@cref}{{[subappendix][3][2147483647,1]A.3}{[1][34][]34}}
\newlabel{sol:eje:3-softmax}{{A.3}{35}{Tema 3}{section*.18}{}}
\newlabel{sol:eje:3-softmax@cref}{{[subappendix][3][2147483647,1]A.3}{[1][35][]35}}
\newlabel{sol:eje:3-nucleo}{{A.3}{35}{Tema 3}{section*.19}{}}
\newlabel{sol:eje:3-nucleo@cref}{{[subappendix][3][2147483647,1]A.3}{[1][35][]35}}
\newlabel{sol:eje:3-word2vec}{{A.3}{35}{Tema 3}{section*.20}{}}
\newlabel{sol:eje:3-word2vec@cref}{{[subappendix][3][2147483647,1]A.3}{[1][35][]35}}
\newlabel{sol:eje:4-vf}{{A.4}{36}{Tema 4}{section*.21}{}}
\newlabel{sol:eje:4-vf@cref}{{[subappendix][4][2147483647,1]A.4}{[1][36][]36}}
\newlabel{sol:eje:4-taylor}{{A.4}{36}{Tema 4}{section*.22}{}}
\newlabel{sol:eje:4-taylor@cref}{{[subappendix][4][2147483647,1]A.4}{[1][36][]36}}
\newlabel{sol:eje:4-mit}{{A.4}{37}{Tema 4}{section*.23}{}}
\newlabel{sol:eje:4-mit@cref}{{[subappendix][4][2147483647,1]A.4}{[1][37][]37}}
\newlabel{sol:eje:4-tasa}{{A.4}{37}{Tema 4}{section*.24}{}}
\newlabel{sol:eje:4-tasa@cref}{{[subappendix][4][2147483647,1]A.4}{[1][37][]37}}
\newlabel{sol:eje:4-necesaria}{{A.4}{38}{Tema 4}{section*.25}{}}
\newlabel{sol:eje:4-necesaria@cref}{{[subappendix][4][2147483647,1]A.4}{[1][38][]38}}
\newlabel{sol:eje:5-vf}{{A.5}{38}{Tema 5}{section*.26}{}}
\newlabel{sol:eje:5-vf@cref}{{[subappendix][5][2147483647,1]A.5}{[1][38][]38}}
\newlabel{sol:eje:5-hessianos}{{A.5}{38}{Tema 5}{section*.27}{}}
\newlabel{sol:eje:5-hessianos@cref}{{[subappendix][5][2147483647,1]A.5}{[1][38][]38}}
\newlabel{sol:eje:5-boyd}{{A.5}{39}{Tema 5}{section*.28}{}}
\newlabel{sol:eje:5-boyd@cref}{{[subappendix][5][2147483647,1]A.5}{[1][39][]39}}
\newlabel{sol:eje:5-red}{{A.5}{39}{Tema 5}{section*.29}{}}
\newlabel{sol:eje:5-red@cref}{{[subappendix][5][2147483647,1]A.5}{[1][39][]39}}
\newlabel{sol:eje:5-softmax-W}{{A.5}{40}{Tema 5}{section*.30}{}}
\newlabel{sol:eje:5-softmax-W@cref}{{[subappendix][5][2147483647,1]A.5}{[1][40][]40}}
\newlabel{sol:eje:5-unico}{{A.5}{40}{Tema 5}{section*.31}{}}
\newlabel{sol:eje:5-unico@cref}{{[subappendix][5][2147483647,1]A.5}{[1][40][]40}}
\newlabel{sol:eje:6-porque}{{A.6}{41}{Tema 6}{section*.32}{}}
\newlabel{sol:eje:6-porque@cref}{{[subappendix][6][2147483647,1]A.6}{[1][41][]41}}
\newlabel{sol:eje:6-cubo}{{A.6}{41}{Tema 6}{section*.33}{}}
\newlabel{sol:eje:6-cubo@cref}{{[subappendix][6][2147483647,1]A.6}{[1][41][]41}}
\newlabel{sol:eje:6-direccional}{{A.6}{42}{Tema 6}{section*.34}{}}
\newlabel{sol:eje:6-direccional@cref}{{[subappendix][6][2147483647,1]A.6}{[1][42][]42}}
\newlabel{sol:eje:6-cs231n}{{A.6}{42}{Tema 6}{section*.35}{}}
\newlabel{sol:eje:6-cs231n@cref}{{[subappendix][6][2147483647,1]A.6}{[1][42][]42}}
\newlabel{sol:eje:6-fermi}{{A.6}{42}{Tema 6}{section*.36}{}}
\newlabel{sol:eje:6-fermi@cref}{{[subappendix][6][2147483647,1]A.6}{[1][42][]42}}
\newlabel{sol:eje:6-hoptimo}{{A.6}{43}{Tema 6}{section*.37}{}}
\newlabel{sol:eje:6-hoptimo@cref}{{[subappendix][6][2147483647,1]A.6}{[1][43][]43}}
\makeatother

\begin{document}
\frontmatter
\begin{titlepage}
\centering
\vspace*{2.5cm}
{\large\scshape Anatomía de un LLM}\par
\vspace{0.4cm}
{\normalsize Fundamentos matemáticos de los modelos de lenguaje}\par
\vspace{2.2cm}
\rule{\textwidth}{0.8pt}\par\vspace{0.5cm}
{\small\scshape Fase I \textperiodcentered\ Fundamentos matemáticos}\par\vspace{0.3cm}
{\Huge\bfseries Módulo 4\par}\vspace{0.3cm}
{\LARGE Cálculo para el aprendizaje automático\par}
\vspace{0.5cm}\rule{\textwidth}{0.4pt}\par
\vspace{1.6cm}
{\large Resumen del módulo}\par
\vfill
{\small Edición de octubre de 2026}\par
\end{titlepage}

\chapter*{Sobre este resumen}
\addcontentsline{toc}{chapter}{Sobre este resumen}
\markboth{Sobre este resumen}{}

Este documento reúne lo indispensable de los apuntes del Módulo 4: el texto expositivo, todas las definiciones y todos los resultados (teoremas, proposiciones, lemas y corolarios), las tablas y las observaciones de contenido. Se han suprimido las demostraciones, los ejemplos, las lecturas, las figuras, los ejercicios y sus soluciones. La numeración de temas, secciones, enunciados y ecuaciones es la misma que en los apuntes, de modo que cualquier referencia puede buscarse allí; las que remiten a elementos suprimidos (un ejemplo, una figura, una ecuación de una demostración) conservan su número en los apuntes.

\section*{Notación}
\begin{center}
\begin{tabular}{@{}ll@{}}
\toprule
Símbolo & Significado \\
\midrule
$x=(x_1,\dots,x_n)\T\in\R^{n}$ & vector columna; $e_i$ es el $i$-ésimo vector canónico \\
$\uno=(1,\dots,1)\T$ & vector de unos \\
$\mathbb{I}[\cdot]$ & indicador: vale 1 si la condición se cumple y 0 en caso contrario \\
$\partial f/\partial x_i$, $\nabla f(x)$ & derivada parcial y gradiente (vector columna) \\
$D_u f(x)$ & derivada direccional en la dirección unitaria $u$ \\
$J_g(x)\in\R^{m\times n}$ & jacobiano de $g\colon\R^{n}\to\R^{m}$, $(J_g)_{ji}=\partial g_j/\partial x_i$ \\
$H_f(x)$, $\nabla^{2}f(x)$ & hessiano, $(H_f)_{ij}=\partial^{2}f/\partial x_i\partial x_j$ \\
$A\succeq 0$, $A\succ 0$ & $A$ simétrica semidefinida o definida positiva \\
$\lambda_{\min}(A)$, $\lambda_{\max}(A)$ & menor y mayor autovalor de una matriz simétrica \\
$\langle A,B\rangle=\tr(A\T B)$ & producto escalar de Frobenius de matrices \\
$z\in\R^{K}$, $p=\softmax(z)$, $y$ & logits, probabilidades y vector \emph{one-hot} del token correcto \\
$\delta=\partial L/\partial y$ & señal de error que llega a la salida de una capa \\
$\eta>0$ & tasa de aprendizaje \\
$\varepsilon_{\text{máq}}$ & épsilon de la máquina (unidad de redondeo) \\
$o(\cdot)$, $O(\cdot)$ & símbolos de Landau \\
\bottomrule
\end{tabular}
\end{center}


\tableofcontents

\mainmatter
% =====================================================================
\chapter{Derivadas parciales y gradiente}
\label{tema:1}
% =====================================================================

La pérdida de un modelo de lenguaje es una función escalar de muchas variables. Este tema define las derivadas parciales y la diferenciabilidad, introduce el gradiente y la derivada direccional, y demuestra que el gradiente señala la dirección de máximo ascenso, que es el fundamento del descenso de gradiente. Se siguen las clases~7, 10 y~11 del curso 18.02 del MIT~\cite{mit1802}, la §5.2 de~\cite{mml} y la §2.4 de~\cite{d2l}.

\section{Derivadas parciales}

Sea $U\subseteq\R^{n}$ un abierto y $f\colon U\to\R$. Para estudiar cómo varía $f$ cerca de un punto, se empieza por moverse en una sola coordenada.

\begin{definicion}[Derivada parcial]\label{def:parcial}
La \emph{derivada parcial} de $f$ respecto de $x_i$ en el punto $x\in U$ es
\[
\frac{\partial f}{\partial x_i}(x)=\lim_{h\to 0}\frac{f(x+h\,e_i)-f(x)}{h},
\]
si el límite existe. Equivale a derivar $f$ como función de una sola variable $x_i$, tratando las demás como constantes.
\end{definicion}

\addtocounter{definicion}{1}%

\addtocounter{definicion}{1}%

\section{Diferenciabilidad y gradiente}

La existencia de las derivadas parciales no basta para que $f$ se comporte bien cerca de $x$: las parciales solo exploran las direcciones de los ejes. La noción adecuada exige que $f$ admita una buena aproximación lineal en todas las direcciones a la vez~\cite[§5.2]{mml}.

\begin{definicion}[Diferenciabilidad y gradiente]\label{def:diferenciable}
La función $f$ es \emph{diferenciable} en $x\in U$ si existe un vector $g\in\R^{n}$ tal que
\[
f(x+h)=f(x)+g\T h+o(\norm{h})\qquad\text{cuando } h\to 0,
\]
es decir, $\lim_{h\to 0}\bigl(f(x+h)-f(x)-g\T h\bigr)/\norm{h}=0$. El vector $g$ se llama \emph{gradiente} de $f$ en $x$ y se escribe $\nabla f(x)$. La función afín $h\mapsto f(x)+\nabla f(x)\T h$ es la \emph{aproximación lineal} (o de primer orden) de $f$ en $x$.
\end{definicion}

\begin{proposicion}\label{prop:grad-parciales}
Si $f$ es diferenciable en $x$, existen todas las parciales en $x$ y
\[
\nabla f(x)=\Bigl(\frac{\partial f}{\partial x_1}(x),\dots,\frac{\partial f}{\partial x_n}(x)\Bigr)\T.
\]
En particular, el gradiente es único.
\end{proposicion}

El gradiente tiene la misma forma que la variable: si $x\in\R^{n}$, entonces $\nabla f(x)\in\R^{n}$. En el \cref{ej:parciales}, $\nabla f(2,1)=(4,7)\T$. Para un LLM, el gradiente de la pérdida tiene tantas componentes como parámetros.

\begin{teorema}[Condición suficiente de diferenciabilidad]\label{teo:C1}
Si todas las derivadas parciales de $f$ existen en un entorno de $x$ y son continuas en $x$, entonces $f$ es diferenciable en $x$. En particular, toda $f\in C^{1}(U)$ es diferenciable en todo punto de $U$.
\end{teorema}

La demostración, estándar en los cursos de cálculo de varias variables, aplica el teorema del valor medio en cada coordenada por separado y se omite aquí. Las funciones que se usan para construir una red (polinomios, exponenciales, logaritmos, cocientes con denominador no nulo y sus composiciones) son de clase $C^{\infty}$ en su dominio. La excepción habitual es la función ReLU, $\max(0,t)$, que no es derivable en $t=0$; el \cref{tema:6} vuelve sobre ella.

\begin{observacion}[Parciales sin diferenciabilidad]
Sea $f(x,y)=xy/(x^{2}+y^{2})$ si $(x,y)\ne(0,0)$ y $f(0,0)=0$. Como $f$ se anula sobre los dos ejes, $\partial f/\partial x(0,0)=\partial f/\partial y(0,0)=0$. Sin embargo, $f(t,t)=1/2$ para todo $t\ne 0$, de modo que $f$ ni siquiera es continua en el origen y, por tanto, no es diferenciable allí. Las parciales existen, pero no son continuas en $(0,0)$.
\end{observacion}

\section{Derivada direccional}

\begin{definicion}[Derivada direccional]\label{def:direccional}
Sea $u\in\R^{n}$ con $\norm{u}=1$. La \emph{derivada direccional} de $f$ en $x$ en la dirección $u$ es
\[
D_u f(x)=\lim_{t\to 0}\frac{f(x+t\,u)-f(x)}{t}.
\]
Las derivadas parciales son las derivadas direccionales en las direcciones $e_1,\dots,e_n$.
\end{definicion}

\begin{proposicion}\label{prop:direccional}
Si $f$ es diferenciable en $x$, entonces $D_u f(x)=\nabla f(x)\T u$ para todo $u$ unitario.
\end{proposicion}

\addtocounter{definicion}{1}%

\section{La dirección de máximo ascenso}

\begin{teorema}[Máximo ascenso]\label{teo:max-ascenso}
Sea $f$ diferenciable en $x$ con $\nabla f(x)\ne 0$. Para todo $u$ unitario,
\[
-\norm{\nabla f(x)}\le D_u f(x)\le\norm{\nabla f(x)}.
\]
La cota superior se alcanza únicamente en $u^{*}=\nabla f(x)/\norm{\nabla f(x)}$ y la inferior únicamente en $-u^{*}$.
\end{teorema}

Así, $\nabla f(x)$ apunta en la dirección en la que $f$ crece más deprisa, su norma es esa tasa máxima de crecimiento y $-\nabla f(x)$ es la dirección de \emph{máximo descenso}.

\begin{lema}[Derivada a lo largo de una curva]\label{lem:curva}
Sea $r\colon(a,b)\to U$ derivable en $t_0$ y $f$ diferenciable en $r(t_0)$. Entonces $t\mapsto f(r(t))$ es derivable en $t_0$ y
\[
\frac{d}{dt}f(r(t))\Big|_{t=t_0}=\nabla f(r(t_0))\T r'(t_0).
\]
\end{lema}

El \cref{lem:curva} es un caso particular de la regla de la cadena, que se estudia en general en el \cref{tema:2}.

\begin{corolario}[Gradiente y conjuntos de nivel]\label{cor:nivel}
Sea $r$ una curva derivable contenida en un conjunto de nivel de $f$, es decir, $f(r(t))=c$ para todo $t$. Entonces $\nabla f(r(t))\T r'(t)=0$: el gradiente es ortogonal a la recta tangente a la curva de nivel.
\end{corolario}

\addtocounter{figure}{1}%

\section{Descenso de gradiente}

El \cref{teo:max-ascenso} sugiere un método para minimizar $f$: desplazarse repetidamente en la dirección $-\nabla f$.

\begin{definicion}[Descenso de gradiente]\label{def:dg}
Dado un punto inicial $x_0$ y una \emph{tasa de aprendizaje} $\eta>0$, el \emph{descenso de gradiente} genera la sucesión
\[
x_{k+1}=x_k-\eta\,\nabla f(x_k),\qquad k=0,1,2,\dots
\]
\end{definicion}

\begin{proposicion}[Descenso a primer orden]\label{prop:descenso}
Si $f$ es diferenciable en $x$ y $g=\nabla f(x)\ne 0$, existe $\eta_0>0$ tal que $f(x-\eta g)<f(x)$ para todo $\eta\in(0,\eta_0)$.
\end{proposicion}

La proposición garantiza el descenso solo para pasos pequeños. El tamaño máximo admisible depende de la curvatura de $f$ y se determina en el \cref{tema:4}.

\addtocounter{definicion}{1}%

\addtocounter{definicion}{1}%

\begin{observacion}[Relación con la física]
Una partícula en un potencial $U(x)$ sometida a un rozamiento muy fuerte se mueve con velocidad proporcional a la fuerza $-\nabla U$: $\dot x(t)=-\nabla U(x(t))$. Esta ecuación se llama \emph{flujo de gradiente}, y la \cref{def:dg} es su discretización por el método de Euler con paso $\eta$. Añadir un término de ruido da la dinámica de Langevin, que sirve de modelo para el descenso de gradiente estocástico.
\end{observacion}


% =====================================================================
\chapter{Regla de la cadena y jacobiano}
\label{tema:2}
% =====================================================================

Una red neuronal es una composición de muchas funciones sencillas, y su pérdida depende de cada parámetro a través de largas cadenas de cálculos intermedios. Este tema generaliza la derivada a funciones con valores vectoriales mediante el jacobiano, demuestra la regla de la cadena en forma matricial y muestra cómo se organiza el cálculo sobre un grafo computacional, en modo directo o en modo inverso. Se sigue~\cite[§5.2, §5.3 y §5.6]{mml}, la clase~10 de~\cite{mit1802} y~\cite[§6.5]{gbc}.

\section{Funciones vectoriales y jacobiano}

\begin{definicion}[Jacobiano]\label{def:jacobiano}
Sea $U\subseteq\R^{n}$ abierto y $g=(g_1,\dots,g_m)\T\colon U\to\R^{m}$. La función $g$ es \emph{diferenciable} en $x$ si existe una matriz $J\in\R^{m\times n}$ tal que
\[
g(x+h)=g(x)+J\,h+o(\norm{h})\qquad\text{cuando } h\to 0.
\]
La matriz $J$ se llama \emph{jacobiano} de $g$ en $x$ y se escribe $J_g(x)$. Aplicando la \cref{prop:grad-parciales} a cada componente, $g$ es diferenciable si y solo si lo es cada $g_j$, y
\[
\bigl(J_g(x)\bigr)_{ji}=\frac{\partial g_j}{\partial x_i}(x),\qquad j=1,\dots,m,\quad i=1,\dots,n.
\]
La fila $j$ de $J_g(x)$ es $\nabla g_j(x)\T$; la columna $i$ es la derivada de $g$ respecto de $x_i$.
\end{definicion}

El jacobiano tiene una fila por cada salida y una columna por cada entrada. Para $m=1$ es una matriz fila y $\nabla f(x)=J_f(x)\T$.

\addtocounter{definicion}{1}%

\addtocounter{definicion}{1}%

\section{La regla de la cadena}

\begin{teorema}[Regla de la cadena]\label{teo:cadena}
Sean $g\colon U\subseteq\R^{n}\to\R^{m}$ diferenciable en $x$ y $f\colon V\subseteq\R^{m}\to\R^{k}$ diferenciable en $g(x)$, con $g(U)\subseteq V$. Entonces $f\circ g$ es diferenciable en $x$ y
\[
J_{f\circ g}(x)=J_f\bigl(g(x)\bigr)\,J_g(x).
\]
\end{teorema}

\begin{corolario}[Caso escalar]\label{cor:cadena-escalar}
Si $f\colon\R^{m}\to\R$ y $u=g(x)$, entonces
\[
\frac{\partial (f\circ g)}{\partial x_i}(x)=\sum_{j=1}^{m}\frac{\partial f}{\partial u_j}(u)\,\frac{\partial u_j}{\partial x_i}(x),
\qquad\text{es decir,}\qquad
\nabla_x (f\circ g)=J_g(x)\T\,\nabla_u f(u).
\]
\end{corolario}

La suma del \cref{cor:cadena-escalar} tiene una lectura directa: si $x_i$ influye en $f$ a través de varias variables intermedias $u_j$, su efecto total es la suma de los efectos transmitidos por cada una de ellas.

\addtocounter{definicion}{1}%

\section{Grafos computacionales}

\begin{definicion}[Grafo computacional]\label{def:grafo}
Un \emph{grafo computacional} es un grafo dirigido acíclico cuyos nodos son variables: los nodos sin predecesores son las entradas, y cada uno de los demás nodos $v$ se calcula como una función diferenciable $v=\phi_v(\text{padres}(v))$ de sus predecesores. Si la salida es un escalar $f$, se llama \emph{adjunto} de un nodo $v$ a $\bar v=\partial f/\partial v$, la derivada de $f$ respecto de $v$ considerando todas las variables posteriores como funciones de $v$~\cite[§6.5.1]{gbc}.
\end{definicion}

\begin{proposicion}[Retropropagación]\label{prop:retro}
En un grafo computacional con salida escalar $f$ se cumple $\bar f=1$ y, para todo nodo $v\neq f$,
\[
\bar v=\sum_{w\,:\,v\in\text{padres}(w)}\bar w\;\frac{\partial\phi_w}{\partial v}.
\]
Por tanto, recorriendo los nodos en orden topológico inverso se obtienen los adjuntos de todas las entradas, es decir, el gradiente de $f$.
\end{proposicion}

\addtocounter{figure}{1}%

\section{Modo directo y modo inverso}

Sea una composición $F=f_L\circ\dots\circ f_1$, con $f_\ell\colon\R^{d_{\ell-1}}\to\R^{d_\ell}$. Por el \cref{teo:cadena},
\[
J_F=J_L\,J_{L-1}\cdots J_1,\qquad J_\ell\in\R^{d_\ell\times d_{\ell-1}},
\]
con cada $J_\ell$ evaluado en el valor intermedio correspondiente. El producto de matrices es asociativo, pero el coste de calcularlo depende del orden en que se agrupen los factores~\cite[§5.6]{mml}.

\begin{proposicion}[Coste de los dos modos]\label{prop:modos}
Multiplicar una matriz $a\times b$ por otra $b\times c$ requiere $abc$ productos escalares. En consecuencia:
\begin{enumerate}
  \item en \emph{modo directo}, se calcula $J_F v=J_L(\cdots(J_2(J_1 v)))$ para un vector $v\in\R^{d_0}$ con $\sum_\ell d_\ell d_{\ell-1}$ productos; obtener el jacobiano completo exige repetirlo $d_0$ veces, una por columna;
  \item en \emph{modo inverso}, se calcula $w\T J_F=((w\T J_L)J_{L-1})\cdots J_1$ para un vector $w\in\R^{d_L}$ con el mismo número de productos; obtener el jacobiano completo exige repetirlo $d_L$ veces, una por fila.
\end{enumerate}
Si la salida es escalar ($d_L=1$), el gradiente se obtiene con una sola pasada en modo inverso.
\end{proposicion}

La retropropagación de la \cref{prop:retro} es el modo inverso aplicado a un grafo. Su coste es un múltiplo pequeño del coste de evaluar $f$, independiente del número de entradas; por ello la pérdida de un modelo con miles de millones de parámetros puede derivarse con una sola pasada hacia atrás. El precio es la memoria: el paso hacia atrás necesita los valores intermedios del paso hacia delante. La implementación de este mecanismo (diferenciación automática) es el objeto del módulo~10.

\addtocounter{definicion}{1}%


% =====================================================================
\chapter{Derivadas matriciales: capa lineal y softmax}
\label{tema:3}
% =====================================================================

Casi todo el cómputo del entrenamiento de un modelo de lenguaje consiste en derivar dos tipos de operación: capas lineales y la combinación de softmax con entropía cruzada en la salida. Este tema define el gradiente de una función escalar respecto de una matriz, obtiene los gradientes de la capa lineal, individualmente y por lotes, y deduce el resultado más usado del curso: el gradiente de la entropía cruzada respecto de los logits es $p-y$. Se sigue~\cite[§5.3--§5.5]{mml}, la tarea~2 de~\cite{cs224n} y, como tabla de consulta, \cite{cookbook}.

\section{Gradiente respecto de una matriz}

\begin{definicion}[Gradiente respecto de una matriz]\label{def:grad-matriz}
Sea $L\colon\R^{m\times n}\to\R$. Su \emph{gradiente} en $W$ es la matriz $\nabla_W L=\partial L/\partial W\in\R^{m\times n}$ de entradas
\[
\Bigl(\frac{\partial L}{\partial W}\Bigr)_{ij}=\frac{\partial L}{\partial W_{ij}}.
\]
Equivale a identificar $\R^{m\times n}$ con $\R^{mn}$ y reordenar el gradiente con la forma de $W$.
\end{definicion}

\begin{proposicion}[Aproximación lineal]\label{prop:frobenius}
Si $L$ es diferenciable en $W$, para toda perturbación $E\in\R^{m\times n}$,
\[
L(W+E)=L(W)+\Bigl\langle\frac{\partial L}{\partial W},E\Bigr\rangle+o(\norm{E}),\qquad
\langle A,B\rangle=\sum_{i,j}A_{ij}B_{ij}=\tr(A\T B).
\]
\end{proposicion}

La regla práctica que se deriva es que \emph{el gradiente de un escalar respecto de un parámetro tiene la misma forma que el parámetro}. Comprobar las formas es la primera verificación de cualquier fórmula de este tema.

\section{La capa lineal}

\begin{proposicion}[Gradientes de una capa lineal]\label{prop:capa}
Sean $W\in\R^{m\times n}$, $x\in\R^{n}$, $b\in\R^{m}$, $y=Wx+b$ y $L=\ell(y)$ con $\ell\colon\R^{m}\to\R$ diferenciable. Sea $\delta=\nabla_y\ell\in\R^{m}$ la \emph{señal de error} que llega a la salida. Entonces
\[
\frac{\partial L}{\partial W}=\delta\,x\T,\qquad
\frac{\partial L}{\partial b}=\delta,\qquad
\frac{\partial L}{\partial x}=W\T\delta.
\]
\end{proposicion}

Las formas encajan: $\delta x\T$ es $(m\times 1)(1\times n)=m\times n$, como $W$, y $W\T\delta$ es $(n\times m)(m\times 1)=n\times 1$, como $x$. El vector $W\T\delta$ es la señal de error que la capa transmite a la capa anterior.

\addtocounter{definicion}{1}%

\begin{corolario}[Capa lineal por lotes]\label{cor:lotes}
Sean $X\in\R^{B\times n}$ una matriz cuyas filas son $B$ ejemplos, $W\in\R^{n\times m}$, $Z=XW\in\R^{B\times m}$ y $L=\ell(Z)$ con $\Delta=\partial L/\partial Z\in\R^{B\times m}$. Entonces
\[
\frac{\partial L}{\partial W}=X\T\Delta,\qquad\frac{\partial L}{\partial X}=\Delta\,W\T.
\]
\end{corolario}

El \cref{cor:lotes} es la \cref{prop:capa} escrita con filas en lugar de columnas y sumada sobre los ejemplos del lote: $X\T\Delta=\sum_b x_b\,\delta_b\T$. Es la convención de las bibliotecas de aprendizaje profundo.

\addtocounter{definicion}{1}%

\section{El jacobiano de la softmax}

Se recuerda del módulo~1 que, para $z\in\R^{K}$, $p=\softmax(z)$ tiene componentes $p_i=e^{z_i}/S$ con $S=\sum_k e^{z_k}$.

\begin{proposicion}[Jacobiano de la softmax]\label{prop:jac-softmax}
\[
\frac{\partial p_i}{\partial z_j}=p_i\bigl(\mathbb{I}[i=j]-p_j\bigr),\qquad\text{es decir,}\qquad J_{\softmax}(z)=\diag(p)-pp\T.
\]
Esta matriz es simétrica y cumple $J_{\softmax}(z)\,\uno=0$.
\end{proposicion}

La igualdad $J_{\softmax}\uno=0$ es la versión infinitesimal de la invariancia por traslación de la softmax (módulo~1, tema~4): sumar la misma constante a todos los logits no cambia las probabilidades.

\section{Entropía cruzada y el gradiente \texorpdfstring{$p-y$}{p − y}}

Sea $c$ el índice del token correcto e $y=e_c$ su vector \emph{one-hot}. La pérdida de entropía cruzada es $L(z)=-\ln p_c$.

\begin{teorema}[Gradiente de la entropía cruzada]\label{teo:p-y}
\[
\nabla_z L=p-y.
\]
\end{teorema}

\addtocounter{definicion}{1}%

\begin{observacion}[Interpretación]
Un paso de descenso $z\leftarrow z-\eta(p-y)$ aumenta el logit correcto en $\eta(1-p_c)$ y disminuye cada logit incorrecto en $\eta p_i$: se penaliza más a los competidores más probables. El gradiente es pequeño cuando el modelo acierta con confianza ($p_c\approx 1$) y grande cuando se equivoca. Como la expresión $p-y$ no contiene cocientes por $p_c$, las bibliotecas implementan la softmax y la entropía cruzada como una sola operación, más estable que componer las dos (módulo~11).
\end{observacion}

\begin{corolario}[Capa de salida de un LLM]\label{cor:salida}
Sean $h\in\R^{d}$ el estado oculto de la última posición, $W\in\R^{K\times d}$ la matriz de salida, $z=Wh$ y $L=-\ln\softmax(z)_c$. Entonces
\[
\frac{\partial L}{\partial W}=(p-y)\,h\T,\qquad\frac{\partial L}{\partial h}=W\T(p-y).
\]
\end{corolario}

Con un vocabulario de $K=10^{5}$ tokens, $W$ tiene $10^{5}$ filas; la fila $i$ recibe la corrección $(p_i-y_i)\,h$, de modo que todas las filas se actualizan en cada paso, no solo la del token correcto.


% =====================================================================
\chapter{Taylor, hessiano y puntos críticos}
\label{tema:4}
% =====================================================================

El gradiente describe la inclinación de una función; la curvatura la describe el hessiano. Este tema introduce las derivadas de segundo orden, demuestra el desarrollo de Taylor de segundo orden en varias variables, lo usa para clasificar puntos críticos y, como aplicación central, determina para qué tasas de aprendizaje converge el descenso de gradiente sobre una función cuadrática. Se sigue~\cite[§5.7 y §5.8]{mml}, \cite[§4.2 y §4.3.1]{gbc} y la clase~9 de~\cite{mit1802}.

\section{El hessiano}

\begin{definicion}[Hessiano]\label{def:hessiano}
Sea $f\colon U\subseteq\R^{n}\to\R$ con derivadas parciales segundas. El \emph{hessiano} de $f$ en $x$ es la matriz $H_f(x)=\nabla^{2}f(x)\in\R^{n\times n}$ de entradas
\[
\bigl(H_f(x)\bigr)_{ij}=\frac{\partial^{2}f}{\partial x_i\,\partial x_j}(x).
\]
Equivalentemente, $H_f$ es el jacobiano del gradiente: $H_f=J_{\nabla f}$.
\end{definicion}

\begin{teorema}[Schwarz]\label{teo:schwarz}
Si $f\in C^{2}(U)$, entonces $\partial^{2}f/\partial x_i\partial x_j=\partial^{2}f/\partial x_j\partial x_i$ en todo punto de $U$. En consecuencia, $H_f(x)$ es simétrica.
\end{teorema}

La demostración se encuentra en los textos de análisis de varias variables y se omite. Al ser simétrico, el hessiano es diagonalizable en una base ortonormal con autovalores reales (teorema espectral, módulo~3), y esos autovalores son las curvaturas de $f$ en las direcciones de los autovectores~\cite[§4.3.1]{gbc}.

\addtocounter{definicion}{1}%

\begin{observacion}[Tamaño del hessiano]
Para $N$ parámetros el hessiano tiene $N^{2}$ entradas. Con $N=7\cdot 10^{9}$ serían $4{,}9\cdot 10^{19}$ números, algo imposible de almacenar. En la práctica solo se calculan productos hessiano-vector $H v$, que cuestan un múltiplo pequeño de un gradiente, o aproximaciones diagonales.
\end{observacion}

\section{Desarrollo de Taylor de segundo orden}

\begin{teorema}[Taylor de segundo orden]\label{teo:taylor}
Sea $f\in C^{2}(U)$ y sean $x\in U$ y $h\in\R^{n}$ tales que el segmento $\{x+th:t\in[0,1]\}$ está contenido en $U$. Entonces:
\begin{enumerate}
  \item (forma de Lagrange) existe $\tau\in(0,1)$ tal que
  \[
  f(x+h)=f(x)+\nabla f(x)\T h+\tfrac12\,h\T H_f(x+\tau h)\,h;
  \]
  \item (forma local) cuando $h\to 0$,
  \[
  f(x+h)=f(x)+\nabla f(x)\T h+\tfrac12\,h\T H_f(x)\,h+o(\norm{h}^{2}).
  \]
\end{enumerate}
\end{teorema}

\addtocounter{definicion}{1}%

\section{Clasificación de puntos críticos}

\begin{definicion}[Punto crítico]
Un punto $x^{*}\in U$ es \emph{crítico} (o estacionario) si $\nabla f(x^{*})=0$. Es un \emph{mínimo local} si existe $r>0$ con $f(x)\ge f(x^{*})$ para $\norm{x-x^{*}}<r$, y un \emph{mínimo local estricto} si la desigualdad es estricta para $x\ne x^{*}$. Un punto crítico que no es mínimo ni máximo local es un \emph{punto de silla}.
\end{definicion}

\begin{teorema}[Condición necesaria de primer orden]\label{teo:fermat}
Si $f$ es diferenciable en $x^{*}$ y $x^{*}$ es un mínimo local, entonces $\nabla f(x^{*})=0$.
\end{teorema}

\begin{teorema}[Criterio del hessiano]\label{teo:criterio}
Sea $f\in C^{2}(U)$ y $x^{*}$ un punto crítico.
\begin{enumerate}
  \item Si $H_f(x^{*})\succ 0$ (todos sus autovalores son positivos), $x^{*}$ es un mínimo local estricto.
  \item Si $H_f(x^{*})\prec 0$ (todos negativos), $x^{*}$ es un máximo local estricto.
  \item Si $H_f(x^{*})$ tiene autovalores de ambos signos, $x^{*}$ es un punto de silla.
\end{enumerate}
\end{teorema}

\begin{table}[h]
\centering\small
\caption{Clasificación de un punto crítico según los autovalores de $H_f(x^{*})$.}\label{tab:criterio}
\begin{tabular}{@{}lll@{}}
\toprule
Autovalores de $H_f(x^{*})$ & Tipo de punto & Ejemplo en $(0,0)$ \\
\midrule
todos $>0$ & mínimo local estricto & $x^{2}+y^{2}$ \\
todos $<0$ & máximo local estricto & $-x^{2}-y^{2}$ \\
de ambos signos & punto de silla & $x^{2}-y^{2}$ \\
alguno nulo, el resto del mismo signo & no concluyente & $x^{4}+y^{2}$, $-x^{4}+y^{2}$ \\
\bottomrule
\end{tabular}
\end{table}

\begin{observacion}[Criterio en dos variables]
Para $n=2$, el producto de los autovalores es $\det H$ y su suma es $\tr H$. Por ello, si $\det H<0$ el punto es de silla, y si $\det H>0$ es un mínimo cuando $f_{xx}>0$ y un máximo cuando $f_{xx}<0$. Es la forma en que se presenta el criterio en~\cite[clase~9]{mit1802}.
\end{observacion}

\addtocounter{definicion}{1}%

\section{Curvatura y descenso de gradiente}

El \cref{teo:taylor} permite analizar el descenso de gradiente más allá de la \cref{prop:descenso}. El caso modelo es una función cuadrática, que coincide con su propio desarrollo de Taylor de segundo orden.

\begin{proposicion}[Descenso de gradiente sobre una cuadrática]\label{prop:dg-cuadratica}
Sea $f(x)=\tfrac12\,x\T Ax$ con $A$ simétrica y autovalores $0<\lambda_1\le\dots\le\lambda_n$. La iteración $x_{k+1}=x_k-\eta\nabla f(x_k)$ converge al mínimo $x^{*}=0$ desde todo punto inicial si y solo si
\[
0<\eta<\frac{2}{\lambda_n}=\frac{2}{\lambda_{\max}(A)}.
\]
\end{proposicion}

\addtocounter{definicion}{1}%

\begin{observacion}[Número de condición]
El factor de contracción en la dirección $i$ es $\abs{1-\eta\lambda_i}$. Con la elección óptima $\eta=2/(\lambda_1+\lambda_n)$, el peor factor vale $(\kappa-1)/(\kappa+1)$, donde $\kappa=\lambda_n/\lambda_1$ es el \emph{número de condición} de $A$. Cuando $\kappa$ es grande, la tasa de aprendizaje está limitada por la dirección de mayor curvatura y el progreso en la de menor curvatura es muy lento. Es el fenómeno de mal condicionamiento descrito en~\cite[§4.2 y §4.3.1]{gbc}, y la motivación de métodos como Adam (módulo~7).
\end{observacion}

\begin{proposicion}[Paso óptimo a lo largo del gradiente]\label{prop:paso-optimo}
Sea $g=\nabla f(x)\ne 0$ y $H=H_f(x)$ con $g\T Hg>0$. La aproximación de segundo orden $q(\eta)=f(x)-\eta\,g\T g+\tfrac12\eta^{2}g\T Hg$ de $f(x-\eta g)$ es mínima en
\[
\eta^{*}=\frac{g\T g}{g\T Hg}.
\]
Si $f$ es cuadrática, $\eta^{*}$ minimiza exactamente $f(x-\eta g)$.
\end{proposicion}

\addtocounter{definicion}{1}%

\begin{observacion}[Relación con la física]
Cerca de un mínimo, $f(x^{*}+h)\approx f(x^{*})+\tfrac12h\T Hh$ es la energía potencial de un sistema de muelles con matriz de rigidez $H$. Los autovalores grandes corresponden a muelles rígidos, que oscilan deprisa; un paso de descenso demasiado largo en esa dirección sobrepasa el equilibrio y el sistema discreto diverge, como establece la \cref{prop:dg-cuadratica}.
\end{observacion}


% =====================================================================
\chapter{Convexidad}
\label{tema:5}
% =====================================================================

En una función convexa, todo mínimo local es global, de modo que un método de descenso no puede quedar atrapado en un mínimo malo. Este tema define los conjuntos y las funciones convexas, demuestra las caracterizaciones de primer y segundo orden, establece las operaciones que conservan la convexidad y las aplica a la entropía cruzada. Termina explicando por qué la pérdida de una red profunda no es convexa. Se sigue~\cite[§3.1, §3.2 y §4.2]{bv} y las transparencias correspondientes de~\cite{ee364a}.

\section{Conjuntos y funciones convexas}

\begin{definicion}[Conjunto convexo]
Un conjunto $C\subseteq\R^{n}$ es \emph{convexo} si, para todo $a,b\in C$ y todo $\lambda\in[0,1]$, el punto $\lambda a+(1-\lambda)b$ pertenece a $C$: contiene el segmento que une dos cualesquiera de sus puntos.
\end{definicion}

\begin{definicion}[Función convexa]\label{def:convexa}
Sea $C\subseteq\R^{n}$ convexo. Una función $f\colon C\to\R$ es \emph{convexa} si
\[
f\bigl(\lambda a+(1-\lambda)b\bigr)\le\lambda f(a)+(1-\lambda)f(b)\qquad\text{para todo } a,b\in C,\ \lambda\in[0,1],
\]
y \emph{estrictamente convexa} si la desigualdad es estricta cuando $a\ne b$ y $\lambda\in(0,1)$. Se dice que $f$ es \emph{cóncava} si $-f$ es convexa.
\end{definicion}

Geométricamente, la cuerda que une $(a,f(a))$ con $(b,f(b))$ queda por encima de la gráfica (\cref{fig:convexa}). Son convexas, por ejemplo, las funciones afines $a\T x+b$ (con igualdad en la definición), toda norma (por la desigualdad triangular y la homogeneidad) y $x\mapsto e^{x}$ en $\R$.

\addtocounter{figure}{1}%

\begin{lema}[Restricción a rectas]\label{lem:rectas}
Sea $C$ abierto y convexo. Una función $f\colon C\to\R$ es convexa si y solo si, para todo $x\in C$ y todo $v\in\R^{n}$, la función de una variable $\varphi(t)=f(x+tv)$ es convexa en el intervalo $\{t:x+tv\in C\}$.
\end{lema}

\section{Condiciones de primer y segundo orden}

\begin{teorema}[Condición de primer orden]\label{teo:primer-orden}
Sea $C$ abierto y convexo y $f\colon C\to\R$ diferenciable. Entonces $f$ es convexa si y solo si
\[
f(b)\ge f(a)+\nabla f(a)\T(b-a)\qquad\text{para todo } a,b\in C.
\]
\end{teorema}

La condición de primer orden dice que el plano tangente en cualquier punto es una cota inferior global de $f$: la información local del gradiente proporciona información global.

\begin{teorema}[Condición de segundo orden]\label{teo:segundo-orden}
Sea $C$ abierto y convexo y $f\in C^{2}(C)$. Entonces $f$ es convexa si y solo si $H_f(x)\succeq 0$ para todo $x\in C$. Si $H_f(x)\succ 0$ para todo $x\in C$, entonces $f$ es estrictamente convexa.
\end{teorema}

\addtocounter{definicion}{1}%

\begin{observacion}
El recíproco de la segunda parte del \cref{teo:segundo-orden} es falso: $f(x)=x^{4}$ es estrictamente convexa en $\R$, pero $f''(0)=0$.
\end{observacion}

\section{Mínimos de funciones convexas}

\begin{teorema}[Mínimos locales y globales]\label{teo:local-global}
Sea $f\colon C\to\R$ convexa. Todo mínimo local de $f$ es un mínimo global. Si además $f$ es diferenciable y $C$ es abierto, un punto $x^{*}$ es mínimo global si y solo si $\nabla f(x^{*})=0$.
\end{teorema}

Para funciones no convexas, $\nabla f=0$ solo es una condición necesaria. Para funciones convexas es también suficiente, y además global. Por eso el descenso de gradiente con una tasa adecuada resuelve los problemas convexos, como la regresión lineal y la regresión logística o softmax del módulo~8.

\section{Operaciones que conservan la convexidad}

\begin{proposicion}\label{prop:operaciones}
\begin{enumerate}
  \item Si $f_1,\dots,f_k$ son convexas y $w_1,\dots,w_k\ge 0$, entonces $\sum_i w_if_i$ es convexa.
  \item Si $f\colon\R^{m}\to\R$ es convexa, $A\in\R^{m\times n}$ y $b\in\R^{m}$, entonces $g(x)=f(Ax+b)$ es convexa.
  \item Si $f_1$ y $f_2$ son convexas, $\max(f_1,f_2)$ es convexa.
\end{enumerate}
\end{proposicion}

\begin{proposicion}[Convexidad de log-suma-exp]\label{prop:lse-convexa}
La función $\LSE(z)=\ln\sum_{k=1}^{K}e^{z_k}$ es convexa en $\R^{K}$. Su hessiano es $H=\diag(p)-pp\T$ con $p=\softmax(z)$, y para todo $v\in\R^{K}$
\[
v\T Hv=\sum_i p_iv_i^{2}-\Bigl(\sum_i p_iv_i\Bigr)^{2}=\Var_{i\sim p}[v_i]\ge 0,
\]
con igualdad si y solo si $v$ es proporcional a $\uno$.
\end{proposicion}

La demostración de~\cite[§3.1]{bv} obtiene la misma desigualdad mediante la de Cauchy--Schwarz.

\begin{corolario}[Convexidad de la entropía cruzada]\label{cor:ce-convexa}
La entropía cruzada $L(z)=-z_c+\LSE(z)$ es convexa en los logits $z$. Para un vector $x$ fijo, $W\mapsto -\ln\softmax(Wx)_c$ es convexa en $W$, y la pérdida media de la regresión softmax sobre un conjunto de datos es convexa en sus parámetros.
\end{corolario}

Para $z=(2,1,0)\T$, el hessiano de la entropía cruzada tiene autovalores $0$, $0{,}119$ y $0{,}371$. El autovalor nulo corresponde a la dirección $\uno$, a lo largo de la cual la pérdida es constante.

\section{Las redes profundas no son convexas}

\begin{proposicion}[Simetría y convexidad]\label{prop:simetria}
Sea $L\colon\R^{N}\to\R$ convexa. El conjunto de sus mínimos globales es convexo. En consecuencia, si $\theta_1$ y $\theta_2$ son mínimos globales y $\tfrac12(\theta_1+\theta_2)$ no lo es, $L$ no es convexa.
\end{proposicion}

En una red con una capa oculta, intercambiar dos neuronas ocultas (junto con sus pesos de entrada y de salida) no cambia la función calculada. Si $\theta_1$ es un mínimo, la permutación $\theta_2$ también lo es, pero el punto medio promedia las dos neuronas y, en general, calcula otra función con pérdida mayor. Por la \cref{prop:simetria}, la pérdida no es convexa.

\addtocounter{definicion}{1}%

\addtocounter{figure}{1}%

\addtocounter{definicion}{1}%


% =====================================================================
\chapter{Diferenciación numérica y comprobación de gradientes}
\label{tema:6}
% =====================================================================

Un error en la implementación de un gradiente no suele provocar un fallo visible: el entrenamiento avanza, pero peor de lo que debería. La herramienta estándar para detectarlo es comparar el gradiente analítico con una aproximación por diferencias finitas. Este tema analiza el error de truncamiento de las diferencias finitas con el teorema de Taylor, el error de redondeo de la aritmética de coma flotante y el compromiso entre ambos, y describe el procedimiento de comprobación de gradientes. Se sigue~\cite[§4.1]{gbc} y las notas de~\cite{cs231n}.

\section{Diferencias finitas}

\begin{definicion}[Diferencias finitas]\label{def:diferencias}
Sea $f\colon\R^{n}\to\R$ y $h>0$. Las \emph{diferencias hacia delante} y \emph{centrada} en la coordenada $i$ son
\[
\Delta_i^{+}f(x)=\frac{f(x+h\,e_i)-f(x)}{h},\qquad
\Delta_i^{c}f(x)=\frac{f(x+h\,e_i)-f(x-h\,e_i)}{2h}.
\]
Ambas aproximan $\partial f/\partial x_i(x)$ cuando $h\to 0$.
\end{definicion}

\begin{proposicion}[Error de truncamiento]\label{prop:truncamiento}
Sea $\phi\colon\R\to\R$ y $x\in\R$.
\begin{enumerate}
  \item Si $\phi\in C^{2}$, existe $\xi\in(x,x+h)$ tal que $\dfrac{\phi(x+h)-\phi(x)}{h}=\phi'(x)+\dfrac{h}{2}\,\phi''(\xi)$.
  \item Si $\phi\in C^{3}$, existe $\xi\in(x-h,x+h)$ tal que $\dfrac{\phi(x+h)-\phi(x-h)}{2h}=\phi'(x)+\dfrac{h^{2}}{6}\,\phi'''(\xi)$.
\end{enumerate}
En consecuencia, el error de la diferencia hacia delante es $O(h)$ y el de la centrada es $O(h^{2})$.
\end{proposicion}

Para una función de varias variables se aplica la proposición a $\phi(t)=f(x+t\,e_i)$.

\addtocounter{definicion}{1}%

\begin{table}[h]
\centering\small
\caption{Diferencias finitas para $f(x)=e^{x}$ en $x=0$ (\cref{ej:exp}).}\label{tab:exp}
\begin{tabular}{@{}lllll@{}}
\toprule
$h$ & Hacia delante & Error & Centrada & Error \\
\midrule
$0{,}1$ & $1{,}051709$ & $5{,}17\cdot 10^{-2}$ & $1{,}001668$ & $1{,}67\cdot 10^{-3}$ \\
$0{,}01$ & $1{,}005017$ & $5{,}02\cdot 10^{-3}$ & $1{,}0000167$ & $1{,}67\cdot 10^{-5}$ \\
$0{,}001$ & $1{,}000500$ & $5{,}00\cdot 10^{-4}$ & $1{,}00000017$ & $1{,}67\cdot 10^{-7}$ \\
\bottomrule
\end{tabular}
\end{table}

\section{Error de redondeo y elección del paso}

En aritmética de coma flotante, cada evaluación de $f$ tiene un error relativo de redondeo del orden del \emph{épsilon de la máquina} $\varepsilon_{\text{máq}}$: unos $1{,}1\cdot 10^{-16}$ en doble precisión (fp64), $6{,}0\cdot 10^{-8}$ en simple precisión (fp32) y $3{,}9\cdot 10^{-3}$ en el formato bf16 que se usa para entrenar modelos de lenguaje~\cite[§4.1]{gbc}. Al restar dos valores casi iguales, $f(x+h)$ y $f(x-h)$, los errores absolutos, de tamaño $\varepsilon_{\text{máq}}\abs{f(x)}$, no disminuyen, y al dividir por $2h$ se amplifican.

\begin{proposicion}[Paso óptimo de la diferencia centrada]\label{prop:h-optimo}
Supóngase que cada evaluación de $f$ tiene un error absoluto a lo sumo $\varepsilon\abs{f}$ y que $\abs{\phi'''}\le M$ cerca de $x$. Entonces el error total de la diferencia centrada está acotado por
\[
E(h)=\frac{\varepsilon\abs{f}}{h}+\frac{M h^{2}}{6},
\]
que alcanza su mínimo en
\[
h^{*}=\Bigl(\frac{3\varepsilon\abs{f}}{M}\Bigr)^{1/3},\qquad E(h^{*})=\frac{M}{2}\,(h^{*})^{2}=O\bigl(\varepsilon^{2/3}\bigr).
\]
\end{proposicion}

Con $\abs{f}\approx M\approx 1$ y doble precisión, $h^{*}\approx(3{,}3\cdot 10^{-16})^{1/3}\approx 7\cdot 10^{-6}$ y el error mínimo es del orden de $10^{-11}$. Pasos más pequeños empeoran el resultado (\cref{fig:error-h}). En simple precisión, $h^{*}\approx 6\cdot 10^{-3}$ y el error mínimo es del orden de $10^{-5}$; en bf16 la comprobación carece de sentido. Por ello, los gradientes se comprueban siempre en doble precisión.

\addtocounter{figure}{1}%

\section{Comprobación de gradientes}

\begin{definicion}[Error relativo]\label{def:error-relativo}
Sean $a$ una componente del gradiente analítico y $n$ la correspondiente aproximación numérica, no ambas nulas. Su \emph{error relativo} es
\[
r(a,n)=\frac{\abs{a-n}}{\max(\abs{a},\abs{n})}.
\]
\end{definicion}

El error relativo, a diferencia del absoluto, no depende de la escala del gradiente. Las notas de~\cite{cs231n} proponen como orientación, para comprobaciones en doble precisión: $r>10^{-2}$ indica casi con seguridad un error; $10^{-4}<r<10^{-2}$ es sospechoso, salvo en funciones con puntos no derivables; $r\le 10^{-7}$ indica un gradiente correcto. En redes profundas, el error relativo tiende a crecer con la profundidad.

\begin{proposicion}[Coste de la comprobación]\label{prop:coste-fd}
Calcular el gradiente completo de $f\colon\R^{N}\to\R$ por diferencias centradas requiere $2N$ evaluaciones de $f$. En cambio, para una dirección $v\in\R^{N}$, la derivada direccional $\nabla f(x)\T v$ se aproxima con dos evaluaciones:
\[
\nabla f(x)\T v=\frac{f(x+hv)-f(x-hv)}{2h}+O(h^{2}).
\]
\end{proposicion}

La retropropagación obtiene el gradiente completo con un coste equivalente a unas pocas evaluaciones de $f$ (\cref{prop:modos}). Las diferencias finitas, por tanto, no sirven para entrenar, sino para \emph{verificar}. El procedimiento habitual es el siguiente~\cite{cs231n}:
\begin{enumerate}[label=\arabic*.]
  \item usar doble precisión y un modelo o un lote pequeños, con pocos datos;
  \item comparar componentes elegidas al azar, o derivadas direccionales en direcciones aleatorias, mediante el error relativo;
  \item tomar $h$ del orden de $10^{-5}$ y repetir con $10^{-4}$ y $10^{-6}$ si el resultado es dudoso;
  \item vigilar los puntos no derivables, como $t=0$ en ReLU: si $x-h$ y $x+h$ quedan a distintos lados de uno de ellos, la aproximación numérica no es fiable.
\end{enumerate}

\addtocounter{definicion}{1}%



\backmatter
% =====================================================================
\begin{thebibliography}{99}
\addcontentsline{toc}{chapter}{Bibliografía}
\markboth{Bibliografía}{}
\small

\bibitem{bv}
S.~Boyd y L.~Vandenberghe.
\newblock \emph{Convex Optimization}.
\newblock Cambridge University Press, 2004.
\newblock Cap.~3, «Convex functions», y cap.~4, «Convex optimization problems». Texto del curso EE364a de Stanford.
\newblock \url{https://web.stanford.edu/~boyd/cvxbook/}

\bibitem{ee364a}
S.~Boyd.
\newblock EE364a: Convex Optimization I.
\newblock Stanford University. Transparencias de la clase «Convex functions».
\newblock \url{https://web.stanford.edu/class/ee364a/lectures/functions.pdf}

\bibitem{cs224n}
C.~Manning \emph{et al.}
\newblock CS224N: Natural Language Processing with Deep Learning.
\newblock Stanford University, invierno de 2021. Tarea~2, «word2vec».
\newblock \url{https://web.stanford.edu/class/archive/cs/cs224n/cs224n.1214/assignments/a2.pdf}

\bibitem{cs231n}
A.~Karpathy y equipo docente de CS231n.
\newblock CS231n: Deep Learning for Computer Vision. Notas del curso, «Neural Networks Part 3: Learning and Evaluation», sección «Gradient checks».
\newblock Stanford University.
\newblock \url{https://cs231n.github.io/neural-networks-3/}

\bibitem{dauphin}
Y.~N. Dauphin, R.~Pascanu, C.~Gulcehre, K.~Cho, S.~Ganguli e Y.~Bengio.
\newblock Identifying and attacking the saddle point problem in high-dimensional non-convex optimization.
\newblock En \emph{Advances in Neural Information Processing Systems 27 (NeurIPS)}, 2014.
\newblock \url{https://arxiv.org/abs/1406.2572}

\bibitem{mml}
M.~P. Deisenroth, A.~A. Faisal y C.~S. Ong.
\newblock \emph{Mathematics for Machine Learning}.
\newblock Cambridge University Press, 2020.
\newblock Cap.~5, «Vector Calculus».
\newblock \url{https://mml-book.github.io/}

\bibitem{gbc}
I.~Goodfellow, Y.~Bengio y A.~Courville.
\newblock \emph{Deep Learning}.
\newblock MIT Press, 2016.
\newblock Cap.~4, «Numerical Computation», y §6.5, «Back-Propagation and Other Differentiation Algorithms».
\newblock \url{https://www.deeplearningbook.org/}

\bibitem{mit1802}
D.~Auroux.
\newblock 18.02: Multivariable Calculus.
\newblock Massachusetts Institute of Technology, MIT OpenCourseWare, primavera de 2006.
\newblock \url{https://ocw.mit.edu/courses/18-02-multivariable-calculus-spring-2006/}

\bibitem{cookbook}
K.~B. Petersen y M.~S. Pedersen.
\newblock \emph{The Matrix Cookbook}.
\newblock Technical University of Denmark, versión del 15 de noviembre de 2012.

\bibitem{d2l}
A.~Zhang, Z.~C. Lipton, M.~Li y A.~J. Smola.
\newblock \emph{Dive into Deep Learning}, versión 1.0.3.
\newblock Cambridge University Press, 2023.
\newblock §2.4, «Calculus».
\newblock \url{https://d2l.ai/}

\end{thebibliography}


\end{document}
